\section{Supplementary proofs and explicit constants}
\label{app:proofs}

\subsection{Restricted scalar nonclosure}
\label{app:scalar-no-go}

We prove Proposition~\ref{prop:scalar-no-go} in the contraction class stated
there.  Consider four trainable deep-linear factors $(u,B,R,v)$ and output
\begin{equation}
 f_n=v^\top RBu.
\end{equation}
Width normalizations do not affect the graph argument; they can be absorbed
into width-dependent scalar coefficients.  A complete contraction is encoded
by a typed graph whose tensor slots are paired and summed.  Multiplication of
contractions is disjoint union of graphs.

\begin{lemma}[Stable contraction independence]
Every fixed finite family of pairwise nonisomorphic typed complete-contraction
graphs gives linearly independent contraction polynomials for every
sufficiently large width.
\end{lemma}

\begin{proof}
Treat tensor entries as independent commuting indeterminates.  For every graph
$G$, partition all index assignments according to which index vertices are
equal.  This writes its contraction as
\begin{equation}
 P_G^{(n)}=\sum_\pi I_{G/\pi}^{(n)},
\end{equation}
where $I_H^{(n)}$ sums only injective labelings of the quotient incidence
network $H$.  Once width exceeds the number of vertices of each type, an
injective labeling produces a tensor-entry monomial that records the typed
incidence network.  It occurs in another injective contraction only for an
isomorphic network.  Hence distinct injective contractions are independent.
In a purported relation among the $P_G^{(n)}$, choose a graph with a maximal
number of index vertices.  Its discrete-partition injective term cannot be
cancelled by any proper quotient, which has fewer vertices, or by a
nonisomorphic graph of the same size.  Descending in vertex count proves that
all coefficients vanish.
\end{proof}

Let $A=RB$ and $H=A^\top A$.  Retain only differentiation of the endpoint
factors in the feature derivative:
\begin{equation}
 \mathcal D_{0,n}v=Au,qquad
 \mathcal D_{0,n}u=A^\top v,qquad
 \mathcal D_{0,n}A=0.
\end{equation}
For $j\ge0$ put
\begin{equation}
 U_j=v^\top AH^ju,quad
 X_j=u^\top H^{j+1}u,quad
 Y_j=v^\top(AA^\top)^{j+1}v.
\end{equation}
Direct differentiation gives
\begin{equation}
 \mathcal D_{0,n}U_j=X_j+Y_j,qquad
 \mathcal D_{0,n}X_j=\mathcal D_{0,n}Y_j=2U_{j+1}.
\end{equation}
Therefore, for $j\ge1$,
\begin{equation}
 \mathcal D_{0,n}^{,2j-1}f_n
 =4^{j-1}\left[u^\top H^ju+v^\top(AA^\top)^jv\right].
 \label{eq:no-go-derivatives}
\end{equation}
In the full feature derivative, every tensor replacement and product-rule
coefficient is nonnegative.  Hence the connected contraction
\begin{equation}
 \Gamma_j=u^\top(B^\top R^\top RB)^ju
\end{equation}
appears with positive coefficient in
$\mathcal D_n^{,2j-1}f_n$.  Its tensor degree is $4j+2$, so the connected
graph types $\Gamma_j$ have unbounded complexity.

Now suppose a fixed encoder uses contractions from a finite bounded-degree
alphabet $\mathcal S$ and evolves by a polynomial autonomous ODE with a
polynomial output readout.  Every time derivative of the output obtained from
that ODE is a polynomial in encoder coordinates.  Its contraction graph is
therefore a disjoint union of connected components from $\mathcal S$.  Choose
$j$ so that $\Gamma_j\notin\mathcal S$.  Equation
\eqref{eq:no-go-derivatives} and stable independence contradict the equality
between the true and encoded $(2j-1)$st output derivatives.  An identity on a
nonempty open state set is a polynomial identity, so the same coefficient
comparison applies globally.

For an analytic finite-jet closure near zero, scale every factor by a common
$\lambda$ and compare the homogeneous Taylor coefficient of the degree of
$\Gamma_j$.  At any fixed power of $\lambda$, only finitely many Taylor
monomials contribute, and their connected components still lie in
$\mathcal S$.  The same contradiction follows.

For squared-loss physical flow, the vector field is a scalar multiple of the
feature field,
\begin{equation}
 \mathcal X_n=-2\kappa(f_n-y)\mathcal D_n.
\end{equation}
If $y\ne0$, the lowest-degree part of
$\mathcal X_n^kf_n$ is $(2\kappa y)^k\mathcal D_n^kf_n$.  If $y=0$, positivity
of graph coefficients shows that the relevant derivative contains a positive
multiple of $f_n^k\mathcal D_n^kf_n$.  In either case a connected component
$\Gamma_j$ outside $\mathcal S$ remains, completing the proof.

The bounded-degree, state-universal, and regular-readout assumptions are
essential.  An unrestricted field can encode the entire future output germ,
and the proposition makes no claim about such oracle encodings or about one
special initialized trajectory.

\subsection{Explicit constants for Theorem~\ref{thm:main}}
\label{app:constants}

This subsection records one concrete choice of the finite constants used in
the main proof.  It also makes clear where width, depth, and horizon enter.

Let
\begin{equation}
 R=\norm{\theta_0}+\sqrt{nT}\rho_0+1,qquad
 \cB=\{\theta:\norm\theta\le R\},qquad \rho_0=\rho(0)>0.
\end{equation}
Put $X=\max_a\norm{u_a}$, $M_0=X$, and recursively define
\begin{align}
s_\ell&=\sup_{|z|\le RM_{\ell-1}}|\phi_\ell'(z)|,
&M_\ell&=\sqrt n\sup_{|z|\le RM_{\ell-1}}|\phi_\ell(z)|,\nonumber\\
B_L&=Rs_L,
&B_\ell&=s_\ell R B_{\ell+1}\quad(\ell<L).
\label{eq:explicit-response-bounds}
\end{align}
These give $\norm{h_{\ell,a}}\le M_\ell$ and
$\norm{\delta_{\ell,a}}\le B_\ell$ on $\cB$.  Let
\begin{align}
Y&=\left(m^{-1}\sum_a y_a^2\right)^{1/2},&
q&=RM_L/n+Y,&S&=Tq,&\ell_T&=1+S,\nonumber\\
v_1&=2B_1X,&
v_\ell&=2B_\ell M_{\ell-1}/n\quad(\ell\ge2),&
v_w&=2M_L,\nonumber\\
V&=\left(v_1^2+\sum_{\ell=2}^Lv_\ell^2+v_w^2\right)^{1/2},
\label{eq:explicit-V}
\end{align}
and
\begin{equation}
 c_f=\frac1n\left[(B_1X)^2+
 \sum_{\ell=2}^L(B_\ell M_{\ell-1})^2+M_L^2\right]^{1/2}.
 \label{eq:explicit-cf}
\end{equation}
Sample Cauchy--Schwarz gives $\norm F\le V\rho$ and the gradient of
each output is bounded by $c_f$.  A finite Lipschitz constant $K$ for $F$ on
the compact convex ball follows from the local Lipschitz activation bounds.

For Variant A, a concrete order-independent derivative recursion is
\begin{align}
\overline Z_1&=S(s_1v_1X)^2,\nonumber\\
\overline Z_\ell&=3s_\ell^2\left[
 S v_\ell^2M_{\ell-1}^2+
 \left(R^2+\frac{2B_\ell^2\ell_T^2M_{\ell-1}^2}{n^2}\right)
 \overline Z_{\ell-1}\right].
\label{eq:explicit-Z}
\end{align}
The derivation is exactly \eqref{eq:A-e}--\eqref{eq:A-chain}, using
$\norm{x+y+z}^2\le3(\norm x^2+\norm y^2+\norm z^2)$.  One may take
\begin{equation}
 \Gamma_A=\frac{\ell_T}{n}\sum_{\ell=2}^L
 B_\ell\sqrt{S\overline Z_{\ell-1}},qquad
 C_A(T)=\Gamma_Ae^{KT}.
 \label{eq:explicit-GammaA}
\end{equation}

For Variant B, define
\begin{equation}
 Q=m\sum_{\ell=2}^L(M_{\ell-1}^2+B_\ell^2),qquad
 C_0=\frac{\ell_TQ}{nm}.
\end{equation}
On the stopped region $\rho\ge\mu/2$, let $J$ be a bound for the derivative
of the current response map $\theta\mapsto\Psi(\theta)$.  Then
\begin{equation}
 \Lambda_T=\ell_T+J(VS+C_0),qquad
 \Gamma_B=\frac{\Lambda_T^2(\Lambda_T-1)}{4nm},qquad
 C_B(T)=\Gamma_Be^{KT}.
 \label{eq:explicit-GammaB}
\end{equation}
Choose $P_A(T)$ so
$\Gamma_Ae^{KT}/\sqrt{P(P+1)}\le1/2$.  Choose $P_B(T)$ so both
\begin{equation}
 \frac{\Gamma_Be^{KT}}{P(P+1)}\le\frac12,
 \qquad
 c_f\frac{\Gamma_Be^{KT}}{P(P+1)}\le\frac\mu4.
\end{equation}
These choices close the first-exit arguments.

\subsection{All orders for bounded activations}

Suppose $\norm{\phi_\ell}_\infty$ and
$\norm{\phi_\ell'}_\infty$ are finite.  Put
$M_\ell=\sqrt n\norm{\phi_\ell}_\infty$ and
$s_\ell=\norm{\phi_\ell'}_\infty$.  The readout equation gives the exact
estimate
\begin{equation}
 \frac\od{\od t}\norm w^2
 =-\frac{4n}{m}\sum_a(f_a-y_a)f_a
 \le nY^2.
\end{equation}
Thus $B_w=(\norm{w_0}^2+nY^2T)^{1/2}$ bounds the readout.  Let
\begin{equation}
 q=B_wM_L/n+Y,qquad S=Tq,qquad \ell_T=1+S,qquad B_L=B_ws_L.
\end{equation}
Proceed downward from layer $L$ to layer two:
\begin{equation}
 D_\ell=\norm{W_{\ell,0}}_F+
 \frac{2\ell_TB_\ell M_{\ell-1}}n,qquad
 B_{\ell-1}=s_{\ell-1}D_\ell B_\ell.
 \label{eq:bounded-downward}
\end{equation}
Projection contraction bounds the Variant A learned increment by the second
term in $D_\ell$; the dense increment is no larger.  Hence
$\norm{\what W_\ell}_F\le D_\ell$, and the backward recurrence gives the next
$B_{\ell-1}$.  Finally,
\begin{equation}
 \norm{W_1(t)}_F\le\norm{W_{1,0}}_F+2SB_1X.
\end{equation}
All physical bounds are independent of $P$.  At fixed order, the moment
integrals are bounded because $1\le\tau\le\ell_T$ and every history is
bounded.  Standard ODE continuation gives global physical-time existence for
every $P$.  Repeating \eqref{eq:explicit-Z} with $D_\ell$ in place of $R$
gives \eqref{eq:old-rate} for all orders.

